\section{Computer-Assisted Verification}\label{sec:certified-package}

The computer-assisted part of the proof consists of two independent
verification steps, each producing a machine-readable certificate.

\textbf{The first step} concerns the geometric realization. Starting from the
line data of a candidate arrangement and the target $\mathrm{O}$-matrix~\eqref{eq:seed19-omatrix},
the normalization script reconstructs the normalized affine
arrangement~\eqref{eq:seed19-normalized-family}
\[
  Y_0:\ y=0,
  \qquad
  L_i:\ y=m_i(x-a_i),
\]
recovers the exact values of the $a_i$ parameters
($\pm\tan(k\pi/18)$, $\pm\eps$), and verifies that the row-order
inequalities~\eqref{eq:row-order-ineqs} and slope-order
conditions~\eqref{eq:slope-separation} hold throughout the
parameter box \eqref{eq:seed19-parameter-box}. The resulting certificate
records:
\begin{itemize}[label=$\triangleright$]
\item the choice of the distinguished line $Y_0$ and the exact labels of the
  $a_i$ parameters;
\item an interval for the special parameter $\eps$;
\item strict ordering of intersections along each line, with a positive
  safety margin;
\item positive slope-separation for every pair of slopes;
\item the direction ordering~\eqref{eq:slope-separation} of the lines,
  verified over the full slope range;
\item the $\eps=0$ checks: substituting $\eps=0$ and evaluating every
  row-order difference using interval arithmetic over the full slope
  range~\eqref{eq:seed19-parameter-box}, exactly two differences vanish
  (the ordering of $Y_0$ and $L_{15}$ along $L_8$, and of $Y_0$ and $L_8$
  along $L_{15}$; all three lines meet at the origin when $a_8=a_{15}=0$)
  and no difference becomes negative;
\item the verdict that the normalized arrangement realizes the target
  $\mathrm{O}$-matrix.
\end{itemize}

\textbf{The second step} concerns the combinatorial properties of the base
configuration. Using the face criterion of Lemma~\ref{lem:face-criterion},
the structural verifier enumerates bounded faces directly from adjacency
relations in the rows of the $\mathrm{O}$-matrix and confirms:
\begin{itemize}[label=$\triangleright$]
\item $Y_0$ touches exactly $17$ triangles;
\item the arrangement has exactly $107$ bounded triangular faces.
\end{itemize}

This separation of steps is not essential. In principle, the combinatorial structure and
triangle counts could be recovered directly from the line equations in a single
verification step. We instead use the $\mathrm{O}$-matrix as an intermediate
object fixing the combinatorial type, so that the two certificates mirror the
two parts of the proof: realization of the $\mathrm{O}$-matrix by the line
arrangement, and verification of the required properties of the
$\mathrm{O}$-matrix itself.

All input data, verification scripts, output certificates,
and detailed format specifications are available at our GitHub repository
\cite{repo}.
